%%%%%%%%%%%%%%%%%%%%%%%%%%%%%%%%%%%%%%%%%%%%%%%%%%%%%%%%%%%%%%%%%%%%%%%%%%%%

%  The following sections are revisions or additions to report.cls
%
%%%%%%%%%%%%%%%%%%%%%%%%%%%%%%%%%%%%%%%%%%%%%%%%%%%%%%%%%%%%%%%%%%%%%%%%%%%%
%  FOOTNOTES:   make them continuously numbered throughout document
%%%%%%%%%%%%%%%%%%%%%%%%%%%%%%%%%%%%%%%%%%%%%%%%%%%%%%%%%%%%%%%%%%%%%%%%%%%%
% define command that can undo having footnotes reset with each chapter:
% (taken from removefr.sty by Donald Arseneau)
\def\@removefromreset#1#2{\let\@tempb\@elt
   \expandafter\let\expandafter\@tempa\csname c@#1\endcsname
   \def\@elt##1{\expandafter\ifx\csname c@##1\endcsname\@tempa\else
         \noexpand\@elt{##1}\fi}%
   \expandafter\edef\csname cl@#2\endcsname{\csname cl@#2\endcsname}%
   \let\@elt\@tempb}
% use the command \@removefromreset to undo the \@addtoreset in report.cls:
\@removefromreset{footnote}{chapter}
%
% define command to allow people to reset footnote counter at will:
\def\resetfootnote{\setcounter{footnote}{0}}  % definition to reset footnote


%%%%%%%%%%%%%%%%%%%%%%%%%%%%%%%%%%%%%%%%%%%%%%%%%%%%%%%%%%%%%%%%%%%%%%%%%
% PAGE LAYOUT
%%%%%%%%%%%%%%%%%%%%%%%%%%%%%%%%%%%%%%%%%%%%%%%%%%%%%%%%%%%%%%%%%%%%%%%%%
% SIDE MARGINS:
%\if@twoside                 % Values for two-sided printing:
%   \oddsidemargin .55in     %   Left margin on odd-numbered pages.
%   \evensidemargin .05in    %   Left margin on even-numbered pages.
%   \marginparwidth 40pt     %   Width of marginal notes.
%\else                       % Values for one-sided printing:
%   \oddsidemargin 0.55in    %   Note that \oddsidemargin = \evensidemargin
%   \evensidemargin 0.55in
%   \marginparwidth 40pt
%\fi
%\marginparsep 10pt          % Horizontal space between outer margin and
%                            % marginal note
%\textwidth 5.9in            % width of text
%
%% VERTICAL SPACING:
%                         % Top of page:
%\topmargin -.5in         %    distance from top of page to running head
%\headheight 14pt         %    Height of box containing running head.
%\headsep .4in            %    Space between running head and text.
%\textheight 8.8in        %    space for text
%\footskip 30pt           %    Distance from baseline of box containing foot
%                         %    to baseline of last line of text.
\RequirePackage[left=1in, right=1in, top=1in, bottom=1in]{geometry}

%%%%%%%%%%%%%%%%%%%%%%%%%%%%%%%%%%%%%%%%%%%%%%%%%%%%%%%%%%%%%%%%%%%%%%%%%%%%
%                            SECTION HEADINGS
%%%%%%%%%%%%%%%%%%%%%%%%%%%%%%%%%%%%%%%%%%%%%%%%%%%%%%%%%%%%%%%%%%%%%%%%%%%%
\newcommand\chaptersize{\large}
\newcommand\sectionsize{\large}
\newcommand\subsectionsize{\normalsize}
\newcommand\subsubsectionsize{\normalsize}
\newcounter{firstchapter}
\setcounter{firstchapter}{0}

\setcounter{secnumdepth}{3}    % Number subsubsections in the chapters
\setcounter{tocdepth}{3}       % Put subsubsections in the table of contents

% Print "CHAPTER" if chap option is specified:
\ifchap
  \renewcommand\@chapapp{\chaptername}
\else
  \renewcommand\@chapapp{}
\fi

\def\specialhead#1{%   GENERAL HEADING WITHOUT A NUMBER (for abstract, etc.)
     \ifx\phantomsection\undefined
     \else
        \clearpage\phantomsection
     \fi
     \addcontentsline{toc}{chapter}{#1}
     \chapter*{\centering #1 \@mkboth{#1}{#1}}}


\def\@chapter[#1]#2{\ifnum\c@firstchapter=0    % start of rpi added stuff
                      \if@twoside\cleardoublepage\suppressfloats[t]\fi
                      \pagenumbering{arabic}
                      \setcounter{firstchapter}{1}
                    \fi
                    \renewcommand{\thepage}{\arabic{page}}
                    \thispagestyle{plain}
                    \pagestyle{myheadings}      % end of rpi added stuff
                    \ifnum \c@secnumdepth >\m@ne
                         \refstepcounter{chapter}%
                         \typeout{\@chapapp\space\thechapter.}%
                         \addcontentsline{toc}{chapter}%
                                   {\protect\numberline{\thechapter.}#1}%
                    \else
                      \addcontentsline{toc}{chapter}{#1}%
                    \fi
                    \chaptermark{#1}%
%                    \addtocontents{lof}{\protect\addvspace{10\p@}}%
%                    \addtocontents{lot}{\protect\addvspace{10\p@}}%
                    \if@twocolumn
                      \@topnewpage[\@makechapterhead{#2}]%
                    \else
                      \@makechapterhead{#2}%
                      \@afterheading
                    \fi}
\def\@makechapterhead#1{%
  \vspace*{0\p@}%
  {\parindent \z@ \raggedright \centering \normalfont \chaptersize
    \ifnum \c@secnumdepth >\m@ne
      \ifchap
         \bfseries \@chapapp{} \thechapter    % print "Chapter" and number
         \vskip -3pt           %\par\nobreak (original)
       \else
         \bfseries \thechapter.
       \fi
    \fi
    \interlinepenalty\@M
    \bfseries #1\par\nobreak
    \vskip 15\p@
  }}

\def\@makeschapterhead#1{%    heading for chapter* command (no numbering)
  \vspace*{0\p@}%
  {\parindent \z@ \raggedright \centering
    \normalfont  \chaptersize
    \interlinepenalty\@M
    \bfseries  #1\par\nobreak
    \vskip 15\p@
  }}

\renewcommand\section{\@startsection {section}{1}{\z@}%
                                   {3.5ex \@plus 1ex \@minus .2ex}%
                                   {.5ex \@plus .3ex}%{1.4ex \@plus.2ex}%
                                   {\normalfont\sectionsize\bfseries}}
\renewcommand\subsection{\@startsection{subsection}{2}{\z@}%
                                     {3.25ex\@plus 1ex \@minus .2ex}%
                                     {.3ex \@plus .2ex}%{1.2ex \@plus .2ex}%
                                     {\normalfont\subsectionsize\bfseries}}
\renewcommand\subsubsection{\@startsection{subsubsection}{3}{\z@}%
                                     {3.25ex\@plus 1ex \@minus .2ex}%
                                     {.2ex \@plus .1ex}%{1ex \@plus .2ex}%
                                     {\normalfont\subsubsectionsize\bfseries}}
% \paragraph and \subparagraph headings unchanged from report.cls.


%%%%%%%%%%%%%%%%%%%%%%%%%%%%%%%%%%%%%%%%%%%%%%%%%%%%%%%%%%%%%%%%%%%%%%%%%%%%
%  APPENDIX
%%%%%%%%%%%%%%%%%%%%%%%%%%%%%%%%%%%%%%%%%%%%%%%%%%%%%%%%%%%%%%%%%%%%%%%%%%%%
\renewcommand\appendix{\par
  \setcounter{chapter}{0}%
  \setcounter{section}{0}%
  \chaptrue
  \renewcommand\@chapapp{\appendixname}%
  \renewcommand\thechapter{\@Alph\c@chapter}}


%%%%%%%%%%%%%%%%%%%%%%%%%%%%%%%%%%%%%%%%%%%%%%%%%%%%%%%%%%%%%%%%%%%%%%%%%%%%
%  FIGURES and TABLES
%%%%%%%%%%%%%%%%%%%%%%%%%%%%%%%%%%%%%%%%%%%%%%%%%%%%%%%%%%%%%%%%%%%%%%%%%%%%
%control float placement:
%
\setcounter{topnumber}{2}
\renewcommand\topfraction{.8}
\setcounter{bottomnumber}{2}
\renewcommand\bottomfraction{.8}
\setcounter{totalnumber}{4}
\renewcommand\textfraction{.2}
\renewcommand\floatpagefraction{.8}
\setcounter{dbltopnumber}{2}
\renewcommand\dbltopfraction{.8}
\renewcommand\dblfloatpagefraction{.8}

\long\def\@makecaption#1#2{%
  \vskip\abovecaptionskip
  \sbox\@tempboxa{\bfseries#1:~~#2}%   Make caption bold
  \ifdim \wd\@tempboxa >\hsize
   {\bfseries #1:~~#2}\par%            Make caption bold
  \else
    \global \@minipagefalse
    \hb@xt@\hsize{\hfil\box\@tempboxa\hfil}%
  \fi
  \vskip\belowcaptionskip}

%%%%%%%%%%%%%%%%%%%%%%%%%%%%%%%%%%%%%%%%%%%%%%%%%%%%%%%%%%%%%%%%%%%%%%%%%%%%
%  TABLE of CONTENTS,  LIST OF TABLES,  LIST OF FIGURES
%%%%%%%%%%%%%%%%%%%%%%%%%%%%%%%%%%%%%%%%%%%%%%%%%%%%%%%%%%%%%%%%%%%%%%%%%%%%

\renewcommand\tableofcontents{%
    \if@twocolumn
      \@restonecoltrue\onecolumn
    \else
      \@restonecolfalse
    \fi
    \chapter*{\contentsname
        \@mkboth{%
           \MakeUppercase\contentsname}{\MakeUppercase\contentsname}}%
    \vskip -1em \begin{singlespace}      % singlespacing
    \@starttoc{toc}%
    \if@restonecol\twocolumn\fi
    \end{singlespace}
    }

\renewcommand*\l@chapter{\pagebreak[3]\vskip 10pt plus 1pt minus 1pt
                         \@dottedtocline{0}{0em}{1.4em}}
\renewcommand*\l@section{\vskip 6pt plus 1pt minus 1pt
                         \@dottedtocline{1}{1.5em}{2.3em}}
\renewcommand*\l@subsection{\ifnum\c@tocdepth>1\vskip 4pt minus 1pt \fi
                         \@dottedtocline{2}{3.8em}{3.2em}}
\renewcommand*\l@subsubsection{\ifnum\c@tocdepth>2 \vskip 3pt minus 1pt \fi
                         \@dottedtocline{3}{7.0em}{4.1em}}

% The following removed because it's not consistent with entries from longtable
%% modify the definition below (taken from latex.ltx) to include
%% "Table" and "Figure" in entries for lot and lof:
%\long\def\@caption#1[#2]#3{%\baselineskip 14.5 pt
%  \addcontentsline{\csname ext@#1\endcsname}{#1}%
%  {\protect\numberline{\csname fnum@#1\endcsname}{\ignorespaces #2}}%
%  \begingroup
%    \@parboxrestore
%    \normalsize
%    \@makecaption{\csname fnum@#1\endcsname}{\ignorespaces #3}\par
%  \endgroup}

\renewcommand\listoffigures{%
    \if@twocolumn
      \@restonecoltrue\onecolumn
    \else
      \@restonecolfalse
    \fi
    \ifx\phantomsection\undefined
      \else
        \clearpage\phantomsection
      \fi
   \chapter*{\centering\listfigurename  % center it
      \@mkboth{\MakeUppercase\listfigurename}%
              {\MakeUppercase\listfigurename}}%
    \addcontentsline{toc}{chapter}{\listfigurename}  % add lof to toc
    \vskip -1em \begin{singlespace}  % singlespacing
    \@starttoc{lof}%
    \if@restonecol\twocolumn\fi
    \end{singlespace}
    }

\renewcommand\listoftables{%
    \if@twocolumn
      \@restonecoltrue\onecolumn
    \else
      \@restonecolfalse
    \fi
    \ifx\phantomsection\undefined
      \else
        \clearpage\phantomsection
      \fi
    \chapter*{\centering\listtablename  % center it
      \@mkboth{\MakeUppercase\listtablename}%
              {\MakeUppercase\listtablename}}%
\addcontentsline{toc}{chapter}{\listtablename}  % add lot to toc
        \vskip -1em \begin{singlespace}  % singlespacing
    \@starttoc{lot}%
    \if@restonecol\twocolumn\fi
    \end{singlespace}
    }

%remove following at same time as remove code to put "Figure" in LOF
%\renewcommand*\l@figure{\vskip 10pt plus 1pt minus 1pt
%                       \@dottedtocline{1}{0em}{5.8em}}

\renewcommand*\l@figure{\vskip 10pt plus 1pt minus 1pt
                       \@dottedtocline{1}{0em}{2.8em}}
\let\l@table\l@figure

\let\rpicaption\caption
\let\lrpicaption\caption


%%%%%%%%%%%%%%%%%%%%%%%%%%%%%%%%%%%%%%%%%%%%%%%%%%%%%%%%%%%%%%%%%%%%%%
%  BIBLIOGRAPHY
%%%%%%%%%%%%%%%%%%%%%%%%%%%%%%%%%%%%%%%%%%%%%%%%%%%%%%%%%%%%%%%%%%%%%%
% RPI def: for use in making an unnumbered bibliography with hanging indents
\def\bibentry{\vskip10pt\par\noindent\hangindent=40pt\frenchspacing}

% report definition modified for no automatic heading and use ragged right
\newcommand{\bibalign}{\raggedright}

\renewenvironment{thebibliography}[1]
     {\bibalign\frenchspacing
      \list{\@biblabel{\@arabic\c@enumiv}}%
           {\settowidth\labelwidth{\@biblabel{#1}}%
            \leftmargin\labelwidth
            \advance\leftmargin\labelsep
            \@openbib@code
            \usecounter{enumiv}%
            \let\p@enumiv\@empty
            \renewcommand\theenumiv{\@arabic\c@enumiv}}%
      \sloppy\clubpenalty4000\widowpenalty4000%
      \sfcode`\.=\@m}
     {\def\@noitemerr
       {\@latex@warning{Empty `thebibliography' environment}}%
      \endlist}


%%%%%%%%%%%%%%%%%%%%%%%%%%%%%%%%%%%%%%%%%%%%%%%%%%%%%%%%%%%%%%%%%%%%%
%  FOR RPI TITLEPAGE, ABSTR. TITLEPAGE & COPYRIGHT PAGE
%%%%%%%%%%%%%%%%%%%%%%%%%%%%%%%%%%%%%%%%%%%%%%%%%%%%%%%%%%%%%%%%%%%%%

 \def\thesistitle#1{\gdef\@thesistitle{#1}}
 \def\author#1{\gdef\@author{#1}}
 \def\degree#1{\gdef\@degree{#1}}
 \def\department#1{\gdef\@department{#1}}
 \newcount\numcomm \numcomm=1
 \def\signaturelines#1{\numcomm=#1}
 \def\thadviser#1{\gdef\@thadviser{#1}}
 \def\projadviser#1{\gdef\@projadviser{#1}} \projadviser{*}
 \def\cothadviser#1{\gdef\@cothadviser{#1}} \cothadviser{*}
 \def\coprojadviser#1{\gdef\@coprojadviser{#1}} \coprojadviser{*}
 \def\cocothadviser#1{\gdef\@cocothadviser{#1}} \cocothadviser{*}
 \def\cocoprojadviser#1{\gdef\@cocoprojadviser{#1}} \cocoprojadviser{*}
 \def\doctype{\if \@projadviser *Dissertation \else Thesis \fi}
 \def\adviser{\if \@projadviser *\@thadviser \else \@projadviser\fi}
 \def\coadviser{\if \@coprojadviser *\@cothadviser\else \@coprojadviser\fi}
 \def\cocoadviser{\if \@cocoprojadviser *\@cocothadviser\else \@cocoprojadviser\fi}
 \def\memberone#1{\gdef\@memberone{#1}}
 \def\membertwo#1{\gdef\@membertwo{#1}}
 \def\memberthree#1{\gdef\@memberthree{#1}}
 \def\memberfour#1{\gdef\@memberfour{#1}}
 \def\memberfive#1{\gdef\@memberfive{#1}}
 \def\membersix#1{\gdef\@membersix{#1}}
 \def\copyrightyear#1{\gdef\@copyrightyear{#1}} \copyrightyear{\the\year}
 \def\submitdate#1{\gdef\@submitdate{#1}}

% Format signature lines on title page
\newcount\numcount
\def\sigblock{
      \numcount=\numcomm
      %\leftline{Approved by the\hfil}
      Approved by:\\
      \adviser, Chair\\
      \advance \numcount by -1
      \if \coadviser * \else % there's a co-adviser
        \coadviser, Co-Chair\\
        \advance \numcount by -1
      \fi
      \ifnum\numcount > 0 %
      \@memberone\\
      \fi
      \advance \numcount by -1 \ifnum\numcount > 0
      \@membertwo\\
      \fi
      \advance \numcount by -1 \ifnum\numcount > 0
      \@memberthree\\
      \fi
      \advance \numcount by -1 \ifnum\numcount > 0
      \@memberfour\\
      \fi
      \advance \numcount by -1 \ifnum\numcount > 0
      \@memberfive\\
      \fi
      \advance \numcount by -1 \ifnum\numcount > 0
      \@membersix\\
      \fi
}

 \def\titlepage{%
     \pagenumbering{roman}
     \thispagestyle{empty}
%     \ifnum\numcomm<4 \vglue .5in\vfil \fi  % masters
%     \ifnum\numcomm>5 \hbox{ } \vspace{-24pt}\fi % make more space on page
     \begin{singlespace}
     \begin{center}
         \parindent=0pt
         {\large\bf\uppercase\expandafter{\@thesistitle}}\\ [20pt]
         {\large\bf{\@author}}\\ [20pt]
         Submitted in Partial Fullfillment of the Requirements\\
         for the Degree of\\ [10pt]
         \emph{\uppercase\expandafter{\@degree}}\\ [40pt]
         \sigblock
%
         \vfill
         \includegraphics[width=0.25\linewidth]{rpi_seal.png}
         \vskip 20pt
         \emph{Department of \expandafter{\@department}}\\
         Rensselaer Polytechnic Institute\\
         Troy, New York
         \vskip 12pt
         \@submitdate
%
%         \ifnum\numcomm < 7 \vskip 32pt \else \vskip 24pt \fi
%         \sigblock
%         \ifnum \numcomm < 7 \vskip 32pt \else \vskip 24pt \fi
%         Rensselaer Polytechnic Institute\\
%         Troy, New York\\ [14pt]
%         \ifnum\numcomm>6 \vskip -8pt \fi
%         \@submitdate
     \end{center}
     \end{singlespace}
     \newpage
     \pagestyle{plain}
%    \pagenumbering{roman}
%    \setcounter{page}{2}
}

\def\copyrightpage{
    \hbox{ }
    \vfill
    \begin{center}
    \copyright\ Copyright \@copyrightyear \\
    by \\
    \@author \\
    All Rights Reserved \\ [12pt]
    \end{center}
    \clearpage}


%%%%%%%%%%%%%%%%%%%%%%%%%%%%%%%%%%%%%%%%%%%%%%%%%%%%%%%%%%%%%%%%%%%%%
% The CODE FROM SETSPACE.STY:
%%%%%%%%%%%%%%%%%%%%%%%%%%%%%%%%%%%%%%%%%%%%%%%%%%%%%%%%%%%%%%%%%%%%%
\def\setspace@size{%
  \ifx \@currsize \normalsize
    \@normalsize
  \else
    \@currsize
  \fi
}

\def\doublespacing{%
  \ifcase \@ptsize \relax % 10pt
    \def \baselinestretch {1.667}%
  \or % 11pt
    \def \baselinestretch {1.618}%
  \or % 12pt
    \def \baselinestretch {1.655}%
  \fi
  \setspace@size
}

\def\onehalfspacing{%
  \ifcase \@ptsize \relax % 10pt
    \def \baselinestretch {1.25}%
  \or % 11pt
    \def \baselinestretch {1.213}%
  \or % 12pt
    \def \baselinestretch {1.241}%
  \fi
  \setspace@size
}

\def\singlespacing{%
  \def \baselinestretch {1}%
  \setspace@size
  \vskip \baselineskip  % Correction for coming into singlespace
}

\def\setstretch#1{%
  \renewcommand{\baselinestretch}{#1}}

%---Stretch the baseline BEFORE calculating the strut size. This improves
%   spacing below tabular environments etc., probably...
%   Comments are welcomed.

\def\@setsize#1#2#3#4{%
  % Modified 1993.04.07--GDG per KPC
  \@nomath#1%
  \let\@currsize#1%
  \baselineskip #2%
  \baselineskip \baselinestretch\baselineskip
  \parskip \baselinestretch\parskip
  \setbox\strutbox \hbox{%
    \vrule height.7\baselineskip depth.3\baselineskip width\z@}%
  \skip\footins \baselinestretch\skip\footins
  \normalbaselineskip\baselineskip#3#4}

%---Increase the space between last line of text and footnote rule.
%\skip\footins 20pt plus4pt minus4pt

%---Reset baselinestretch within floats and footnotes.

% GT:  This is where the conflict with the combination of the color
% package and the figure environment used to occur.

\let\latex@xfloat=\@xfloat
\def\@xfloat #1[#2]{%
  \latex@xfloat #1[#2]%
  \def \baselinestretch {1}\@normalsize  %%  GT: need \@normalsize?
  \normalsize
}

\long\def\@footnotetext#1{%
  \insert\footins{%
    \def\baselinestretch {1}%
    \reset@font\footnotesize
    \interlinepenalty\interfootnotelinepenalty
    \splittopskip\footnotesep
    \splitmaxdepth \dp\strutbox \floatingpenalty \@MM
    \hsize\columnwidth
    \@parboxrestore
    \protected@edef\@currentlabel{%
      \csname p@footnote\endcsname\@thefnmark}%
    \color@begingroup
      \@makefntext{%
        \rule\z@\footnotesep\ignorespaces#1\@finalstrut\strutbox}%
    \color@endgroup}}

\long\def\@mpfootnotetext#1{%
  \global\setbox\@mpfootins\vbox{%
    \unvbox \@mpfootins
    \def\baselinestretch {1}%
    \reset@font\footnotesize
    \hsize\columnwidth
    \@parboxrestore
    \protected@edef\@currentlabel{%
      \csname p@mpfootnote\endcsname\@thefnmark}%
    \color@begingroup
      \@makefntext{%
       \rule\z@\footnotesep\ignorespaces#1\@finalstrut\strutbox}%
   \color@endgroup}}

% A single spaced quote (say) is done by surrounding singlespace with quote.
%
\def\singlespace{%
  \vskip \baselineskip
  \def\baselinestretch {1}%
  \setspace@size
  \vskip -\baselineskip
}

\def\endsinglespace{\par}

%  spacing, doublespace and onehalfspace all are meant to INCREASE the
%  spacing (i.e. calling onehalfspace from within doublespace will not
%  produce a graceful transition between spacings)
%
% Next two definitions fixed for consistency with TeX 3.x

\def\spacing#1{%
  \par
  \begingroup             % moved from \endspacing by PGBR 29-1-91
    \def\baselinestretch {#1}%
    \setspace@size
}

\def\endspacing{%
    \par
    \vskip \parskip
    \vskip \baselineskip
  \endgroup
  \vskip -\parskip
  \vskip -\baselineskip
}

% one and a half spacing is 1.5 x pt size
\def\onehalfspace{%
  \ifcase \@ptsize \relax  % 10pt
    \spacing{1.25}%
  \or % 11pt
    \spacing{1.213}%
  \or % 12pt
    \spacing{1.241}%
  \fi
}

\let\endonehalfspace=\endspacing

% double spacing is 2 x pt size
\def\doublespace{%
  \ifcase \@ptsize \relax % 10pt
    \spacing{1.667}%
  \or % 11pt
    \spacing{1.618}%
  \or % 12pt
    \spacing{1.655}%
  \fi
}

\let\enddoublespace=\endspacing

% GT:  EMSH chose to omit display math part that follows.
% She wrote (see above) that the "altered spacing before and after displayed
% equations ... just looked too much".
%
% Fix up spacing before and after displayed math
% (arraystretch seems to do a fine job for inside LaTeX displayed math,
% since array and eqnarray seem to be affected as expected).
% Changing \baselinestretch and doing a font change also works if done here,
% but then you have to change @setsize to remove the call to @nomath)
%
\everydisplay\expandafter{%
  \the\everydisplay
  \abovedisplayskip \baselinestretch\abovedisplayskip
  \belowdisplayskip \abovedisplayskip
  \abovedisplayshortskip \baselinestretch\abovedisplayshortskip
  \belowdisplayshortskip \baselinestretch\belowdisplayshortskip
}
%%%%%%%%%%%%%%%%%%%%%%  End of setspace.sty %%%%%%%%%%%%%%%%%%%%%%%%%%%%

%% Setup captions
%\RequirePackage[justification=centering]{caption}
%\RequirePackage[justification=centerlast]{caption}
%\RequirePackage[format=hang,justification=centerlast]{caption}
\RequirePackage[format=hang,justification=raggedright,singlelinecheck=off]{caption}
\renewcommand{\captionfont}{\bfseries}

\usepackage[backend=bibtex, style = ieee]{biblatex}